% Options for packages loaded elsewhere
\PassOptionsToPackage{unicode}{hyperref}
\PassOptionsToPackage{hyphens}{url}
\PassOptionsToPackage{dvipsnames,svgnames,x11names}{xcolor}
\documentclass[
  11pt,
  UTF8,
  a4paper,
  fontset=fandol]{ctexart}
\usepackage{xcolor}
\usepackage[margin=2.3cm]{geometry}
\usepackage{amsmath,amssymb}
\setcounter{secnumdepth}{-\maxdimen} % remove section numbering
\usepackage{iftex}
\ifPDFTeX
  \usepackage[T1]{fontenc}
  \usepackage[utf8]{inputenc}
  \usepackage{textcomp} % provide euro and other symbols
\else % if luatex or xetex
  \usepackage{unicode-math} % this also loads fontspec
  \defaultfontfeatures{Scale=MatchLowercase}
  \defaultfontfeatures[\rmfamily]{Ligatures=TeX,Scale=1}
\fi
\usepackage{lmodern}
\ifPDFTeX\else
  % xetex/luatex font selection
  \setmonofont[]{DejaVu Sans Mono}
\fi
% Use upquote if available, for straight quotes in verbatim environments
\IfFileExists{upquote.sty}{\usepackage{upquote}}{}
\IfFileExists{microtype.sty}{% use microtype if available
  \usepackage[]{microtype}
  \UseMicrotypeSet[protrusion]{basicmath} % disable protrusion for tt fonts
}{}
\makeatletter
\@ifundefined{KOMAClassName}{% if non-KOMA class
  \IfFileExists{parskip.sty}{%
    \usepackage{parskip}
  }{% else
    \setlength{\parindent}{0pt}
    \setlength{\parskip}{6pt plus 2pt minus 1pt}}
}{% if KOMA class
  \KOMAoptions{parskip=half}}
\makeatother
\usepackage{listings}
\newcommand{\passthrough}[1]{#1}
\lstset{defaultdialect=[5.3]Lua}
\lstset{defaultdialect=[x86masm]Assembler}
\usepackage{longtable,booktabs,array}
\usepackage{calc} % for calculating minipage widths
% Correct order of tables after \paragraph or \subparagraph
\usepackage{etoolbox}
\makeatletter
\patchcmd\longtable{\par}{\if@noskipsec\mbox{}\fi\par}{}{}
\makeatother
% Allow footnotes in longtable head/foot
\IfFileExists{footnotehyper.sty}{\usepackage{footnotehyper}}{\usepackage{footnote}}
\makesavenoteenv{longtable}
\setlength{\emergencystretch}{3em} % prevent overfull lines
\providecommand{\tightlist}{%
  \setlength{\itemsep}{0pt}\setlength{\parskip}{0pt}}
\lstset{
  basicstyle=\ttfamily\small,
  commentstyle=\upshape,
  breaklines=true,
  breakatwhitespace=false,
  columns=fullflexible,
  keepspaces=true,
  showstringspaces=false
}
\usepackage{bookmark}
\IfFileExists{xurl.sty}{\usepackage{xurl}}{} % add URL line breaks if available
\urlstyle{same}
\hypersetup{
  pdftitle={FreeBSD 构建、验证与 MM/Jail 实验手册},
  colorlinks=true,
  linkcolor={NavyBlue},
  filecolor={Maroon},
  citecolor={Blue},
  urlcolor={RoyalBlue},
  pdfcreator={LaTeX via pandoc}}

\title{FreeBSD 构建、验证与 MM/Jail 实验手册}
\author{}
\date{2026-11-22}

\begin{document}
\maketitle

{
\hypersetup{linkcolor=}
\setcounter{tocdepth}{3}
\tableofcontents
}
\begin{quote}
从交叉构建开始，经过 QEMU、ATF/Kyua 与 GDB 验证，再评估 NixBSD
和容器化开发环境，最后回到更适合内核研究的 KVM+QEMU+GDB 主线。

本文区分五种状态：\textbf{已实测}、\textbf{仅源码审查}、\textbf{未验证}、\textbf{暂停}与\textbf{未开始}。命令中的路径变量需要读者按自己的工作区设置，不代表某台真实机器上的目录。
\end{quote}

\section{0.build on nixos}\label{build-on-nixos}

\subsection{0.1
先划清宿主与目标的边界}\label{ux5148ux5212ux6e05ux5bbfux4e3bux4e0eux76eeux6807ux7684ux8fb9ux754c}

在非 FreeBSD 系统上构建 FreeBSD，最重要的不是把宿主机变成
FreeBSD，而是把宿主工具、目标文件和运行环境分开。

本文的已验证路径是：

\begin{lstlisting}
NixOS host
└── Ubuntu 24.04 build container
    ├── Clang/LLD 20
    ├── FreeBSD source
    ├── external object directory
    └── world + GENERIC kernel + release image
         ↓
      QEMU guest
\end{lstlisting}

NixOS 在这里提供容器、QEMU 和调试工具。\textbf{成功的是 NixOS 宿主上的
Ubuntu 容器交叉构建，不是原生 NixOS 构建。} 原生路径曾遇到非 FHS
布局问题，但一次失败不足以证明它永远不可行，因此本文不提供``原生 NixOS
已成功''的复现承诺。

整个流程只围绕三个问题展开：

\begin{enumerate}
\def\labelenumi{\arabic{enumi}.}
\tightlist
\item
  如何得到一组可信的 world 和 kernel；
\item
  如何把安装结果变成可检查、可启动的磁盘；
\item
  如何保留证据，而不是只相信构建命令的退出码。
\end{enumerate}

\subsection{0.2
工作区是一组约定，不是一份环境清单}\label{ux5de5ux4f5cux533aux662fux4e00ux7ec4ux7ea6ux5b9aux4e0dux662fux4e00ux4efdux73afux5883ux6e05ux5355}

下面的变量贯穿本章：

\begin{lstlisting}[language=sh]
export WORK="$PWD"
export SRC="$WORK/src"
export OBJ="$WORK/obj"
export IMAGES="$WORK/images"
export BUILDER=freebsd-src-builder:ubuntu24.04
export JOBS=4

mkdir -p "$OBJ" "$IMAGES"
\end{lstlisting}

建议让源码、对象和镜像彼此独立：

\begin{lstlisting}
WORK/
├── src/       # FreeBSD source
├── obj/       # bmake、world、kernel 和 release 对象
├── images/    # raw 与 qcow2
└── docker/    # builder 定义和 VM release 配置
\end{lstlisting}

外部对象目录的价值不仅是保持源码树整洁。它还让干净构建、增量构建和不同工具链之间的隔离变得可操作。不要用来源不明的对象目录拼接一个新结果。

获取并固定源码：

\begin{lstlisting}[language=sh]
git clone https://github.com/freebsd/freebsd-src.git "$SRC"
git -C "$SRC" checkout --detach \
    e6baf4c7430d9def848f70e6be7934e8c3df15a0

git -C "$SRC" status --short
git -C "$SRC" log -1 --oneline
\end{lstlisting}

固定提交使``源码是否变化''成为一个明确问题。已有工作区不要直接覆盖
clone，也不要为了复现而丢弃本地修改。

\subsection{0.3 工具链：真正需要解释的是
Clang}\label{ux5de5ux5177ux94feux771fux6b63ux9700ux8981ux89e3ux91caux7684ux662f-clang}

交叉构建会涉及许多宿主工具，但多数只是构建器的实现细节。对这个源码快照而言，值得单独讨论的是
\textbf{Clang/LLD 的版本}。

\passthrough{\lstinline!\_\_is\_nothrow\_convertible!}
内建支持问题出现在较旧工具链的 libc++ C++23 路径中。已实测的完整构建改用
Clang/LLD 20 后，world 和 kernel 均完成构建。这不是通过修改 FreeBSD
源码绕过问题，而是让编译器满足源码当时的需要。

这里有一条比版本号更通用的规则：

\begin{quote}
切换编译器后，不要沿用旧对象目录做``增量验证''。
\end{quote}

旧对象可能隐藏头文件、内建函数或链接行为的变化。首次构建和工具链切换使用
\passthrough{\lstinline!--clean!}；只有相同源码、相同工具链和相同配置下才使用
\passthrough{\lstinline!--no-clean!}。

构建镜像后，先检查真正会被调用的工具：

\begin{lstlisting}[language=sh]
docker build --pull=false -t "$BUILDER" "$WORK/docker"

docker run --rm "$BUILDER" bash -c '
    clang-20 --version
    ld.lld-20 --version
    bmake -V MAKE_VERSION
    python3 --version
'
\end{lstlisting}

版本检查不是装饰。它防止镜像标签与实际工具链脱节。

\subsection{0.4 world 与 kernel}\label{world-ux4e0e-kernel}

FreeBSD 为交叉构建提供
\passthrough{\lstinline!tools/build/make.py!}。它负责自举
bmake、设置交叉工具链并建立目标构建环境，不应简单替换为宿主 GNU make。

首次构建：

\begin{lstlisting}[language=sh]
docker run --rm --name freebsd-src-build \
    --user "$(id -u):$(id -g)" \
    -v "$SRC:/work/src" \
    -v "$OBJ:/work/obj" \
    -w /work/src \
    -e MAKEOBJDIRPREFIX=/work/obj \
    "$BUILDER" \
    ./tools/build/make.py \
        --debug \
        --clean \
        --cross-bindir=/usr/lib/llvm-20/bin \
        TARGET=amd64 \
        TARGET_ARCH=amd64 \
        KERNCONF=GENERIC \
        -DWITH_DISK_IMAGE_TOOLS_BOOTSTRAP \
        buildworld buildkernel \
        -s -j"$JOBS"
\end{lstlisting}

几个参数构成了构建的骨架：

\begin{longtable}[]{@{}
  >{\raggedright\arraybackslash}p{(\linewidth - 2\tabcolsep) * \real{0.4800}}
  >{\raggedright\arraybackslash}p{(\linewidth - 2\tabcolsep) * \real{0.5200}}@{}}
\toprule\noalign{}
\begin{minipage}[b]{\linewidth}\raggedright
参数
\end{minipage} & \begin{minipage}[b]{\linewidth}\raggedright
作用
\end{minipage} \\
\midrule\noalign{}
\endhead
\bottomrule\noalign{}
\endlastfoot
\passthrough{\lstinline!MAKEOBJDIRPREFIX!} & 把对象放到源码树外 \\
\passthrough{\lstinline!TARGET!} /
\passthrough{\lstinline!TARGET\_ARCH!} & 选择目标机器和 ABI \\
\passthrough{\lstinline!KERNCONF!} & 选择内核配置 \\
\passthrough{\lstinline!--cross-bindir!} & 固定交叉 Clang/LLD 入口 \\
\passthrough{\lstinline!-DWITH\_DISK\_IMAGE\_TOOLS\_BOOTSTRAP!} &
自举后续组盘所需工具 \\
\end{longtable}

相同条件下增量重建时，把 \passthrough{\lstinline!--clean!} 改为
\passthrough{\lstinline!--no-clean!}。不要把增量构建当作工具链变化后的干净验证。

\subsection{0.5
从安装树到磁盘镜像}\label{ux4eceux5b89ux88c5ux6811ux5230ux78c1ux76d8ux955cux50cf}

release 流程会再次设置环境。容器内需要显式恢复自举工具路径，避免脚本重置
\passthrough{\lstinline!PATH!} 后找不到刚刚构建的工具：

\begin{lstlisting}[language=sh]
OBJROOT=/work/obj/work/src/amd64.amd64
BUILD_PATH=/work/obj/bmake-install/bin
BUILD_PATH="$BUILD_PATH:$OBJROOT/tmp/bin:$OBJROOT/tmp/usr/sbin"
BUILD_PATH="$BUILD_PATH:$OBJROOT/tmp/usr/bin:$OBJROOT/tmp/legacy/usr/sbin"
BUILD_PATH="$BUILD_PATH:$OBJROOT/tmp/legacy/usr/bin:$OBJROOT/tmp/legacy/bin"
BUILD_PATH="$BUILD_PATH:$OBJROOT/tmp/legacy/usr/libexec"
export FREEBSD_BUILD_PATH="$BUILD_PATH:/usr/local/sbin:/usr/local/bin:/usr/sbin:/usr/bin:/sbin:/bin"
export PATH="$FREEBSD_BUILD_PATH"
\end{lstlisting}

下面的命令在 builder 内执行，调用 release 的
\passthrough{\lstinline!vm-image!} 路径：

\begin{lstlisting}[language=sh]
./tools/build/make.py \
    --no-clean \
    --cross-bindir=/usr/lib/llvm-20/bin \
    TARGET=amd64 TARGET_ARCH=amd64 KERNCONF=GENERIC \
    -DWITH_DISK_IMAGE_TOOLS_BOOTSTRAP \
    -C /work/src/release \
    NOPKGBASE=1 NO_ROOT=1 WITHOUT_QEMU=1 PKG_CMD=true \
    -DWITH_VMIMAGES \
    VMFORMATS=raw VMFSLIST=ufs VMSIZE=5g SWAPSIZE=512m \
    VM_IMAGE_CONFIG=/work/docker/vm-qemu.conf \
    vm-image -s
\end{lstlisting}

\passthrough{\lstinline!NO\_ROOT!} 和 METALOG 让安装树可以在非
root/chroot 条件下生成；\passthrough{\lstinline!WITHOUT\_QEMU!} 避免在
Linux 构建机上执行目标 FreeBSD
二进制。测试镜像的账号和网络配置应只用于隔离环境，不要把空密码 root SSH
当成通用模板。

检查安装树，而不是立即相信镜像：

\begin{lstlisting}[language=sh]
export STAGE="$OBJ/work/src/amd64.amd64/release/vm-image-raw-ufs"

test -x "$STAGE/bin/sh"
test -f "$STAGE/boot/kernel/kernel"
cat "$STAGE/etc/fstab"
\end{lstlisting}

\subsection{0.6
退出码之后，还要检查产物}\label{ux9000ux51faux7801ux4e4bux540eux8fd8ux8981ux68c0ux67e5ux4ea7ux7269}

这个源码快照的 release
路径揭示了一个典型问题：某个组盘函数失败后，上层规则仍可能因
\passthrough{\lstinline!|| true!} 返回成功。于是出现``make 返回
0，却没有新镜像''的假成功。

因此验收至少应包含：

\begin{lstlisting}[language=sh]
qemu-img info "$IMAGES/freebsd-main-amd64.raw"
fdisk -l "$IMAGES/freebsd-main-amd64.raw"
\end{lstlisting}

如果需要手工恢复组盘，应先理解 \passthrough{\lstinline!vmimage.subr!}
的状态和 \passthrough{\lstinline!DESTDIR!}，而不是重复执行整个
release。尤其注意：

\begin{itemize}
\tightlist
\item
  \passthrough{\lstinline!MAKEOBJDIRPREFIX!} 必须与构建时一致；
\item
  \passthrough{\lstinline!vm\_create\_disk!} 会修改
  \passthrough{\lstinline!fstab!}；
\item
  已有 ESP 记录时不要对同一安装树重复组盘；
\item
  恢复失败时保留安装树和日志，避免叠加第二次修改。
\end{itemize}

转换为 qcow2：

\begin{lstlisting}[language=sh]
qemu-img convert -p -f raw -O qcow2 \
    "$IMAGES/freebsd-main-amd64.raw" \
    "$IMAGES/freebsd-main-amd64.qcow2"

qemu-img info "$IMAGES/freebsd-main-amd64.qcow2"
qemu-img check "$IMAGES/freebsd-main-amd64.qcow2"
\end{lstlisting}

同名文件可能已经是有价值的旧镜像，转换前应另行备份。\passthrough{\lstinline!qemu-img check!}
的结果和实际启动是两类证据，不能互相替代。

\subsection{0.7
本章验收边界}\label{ux672cux7ae0ux9a8cux6536ux8fb9ux754c}

\textbf{已实测：} 固定提交的 world、GENERIC kernel、raw/qcow2
生成，以及后续 QEMU 启动。

\textbf{未验证：} 原生 NixOS 构建、KVM 启动和当前快照之外的一般性结论。

构建章最终应交付的不是一串命令，而是一组可以被下一章消费的对象：源码提交、内核符号、安装树、可启动磁盘，以及它们之间明确的对应关系。

\section{1.atf bsd}\label{atf-bsd}

\subsection{1.1
验证不是启动画面之后的一句``能进系统''}\label{ux9a8cux8bc1ux4e0dux662fux542fux52a8ux753bux9762ux4e4bux540eux7684ux4e00ux53e5ux80fdux8fdbux7cfbux7edf}

内核开发需要一套可重复的观察方法。本文把验证环境拆成四个接口：

\begin{longtable}[]{@{}
  >{\raggedright\arraybackslash}p{(\linewidth - 2\tabcolsep) * \real{0.4800}}
  >{\raggedright\arraybackslash}p{(\linewidth - 2\tabcolsep) * \real{0.5200}}@{}}
\toprule\noalign{}
\begin{minipage}[b]{\linewidth}\raggedright
接口
\end{minipage} & \begin{minipage}[b]{\linewidth}\raggedright
职责
\end{minipage} \\
\midrule\noalign{}
\endhead
\bottomrule\noalign{}
\endlastfoot
Serial & 启动日志、登录、ATF/Kyua 和 guest 内状态 \\
QMP & QEMU 运行状态与生命周期 \\
GDB stub & vCPU 寄存器、内核指令和符号 \\
qcow2 snapshot & 把实验写入限制在临时层 \\
\end{longtable}

这四个接口回答的问题不同。串口正常不代表 QMP 可管理，QMP 显示 running
不代表 guest 没有内核错误，GDB 能写寄存器也不代表一次 fuzz 已经完成。

\subsection{1.2 启动一个可丢弃的
VM}\label{ux542fux52a8ux4e00ux4e2aux53efux4e22ux5f03ux7684-vm}

设置验证对象：

\begin{lstlisting}[language=sh]
export IMAGE="$IMAGES/freebsd-main-amd64.qcow2"
export KERNEL_DEBUG="$OBJ/work/src/amd64.amd64/sys/GENERIC/kernel.debug"
export RUNDIR="$(mktemp -d /tmp/freebsd-verify.XXXXXX)"
chmod 700 "$RUNDIR"
\end{lstlisting}

启动：

\begin{lstlisting}[language=sh]
qemu-system-x86_64 \
    -name freebsd-verify \
    -machine q35,accel=tcg \
    -cpu qemu64 \
    -smp 2 -m 2048 \
    -drive if=none,file="$IMAGE",format=qcow2,id=disk0 \
    -device virtio-blk-pci,drive=disk0 \
    -nic user,model=e1000 \
    -boot order=c \
    -display none -monitor none \
    -chardev socket,id=serial0,path="$RUNDIR/serial.sock",server=on,wait=off,logfile="$RUNDIR/serial.log" \
    -serial chardev:serial0 \
    -qmp unix:"$RUNDIR/qmp.sock",server=on,wait=off \
    -gdb unix:"$RUNDIR/gdb.sock",server=on,wait=off \
    -snapshot \
    -daemonize \
    -pidfile "$RUNDIR/qemu.pid" \
    -D "$RUNDIR/qemu.log"
\end{lstlisting}

\passthrough{\lstinline!-snapshot!}
是实验的默认保险。只有明确需要持久修改 guest
时才应移除它，而且每次只能有一个写实例使用基础磁盘。

有权限访问 \passthrough{\lstinline!/dev/kvm!} 的机器可以把
\passthrough{\lstinline!accel=tcg!} 换成
\passthrough{\lstinline!accel=kvm!}，并重新验证全部流程。\textbf{本文的实际寄存器与
ATF 结果来自 TCG，不把 KVM 写成已验证结果。}

\subsection{1.3
QMP：先确认机器仍然受控}\label{qmpux5148ux786eux8ba4ux673aux5668ux4ecdux7136ux53d7ux63a7}

下面的小函数只依赖 Python 标准库，可用于查询和退出：

\begin{lstlisting}[language=sh]
qmp() {
    python3 - "$RUNDIR/qmp.sock" "$@" <<'PY'
import json
import socket
import sys

s = socket.socket(socket.AF_UNIX, socket.SOCK_STREAM)
s.settimeout(10)
s.connect(sys.argv[1])
f = s.makefile("rwb", buffering=0)
print(f.readline().decode().strip())
for name in sys.argv[2:]:
    f.write((json.dumps({"execute": name}) + "\n").encode())
    print(f.readline().decode().strip())
s.close()
PY
}

qmp qmp_capabilities query-status query-version
\end{lstlisting}

健康检查至少要求 \passthrough{\lstinline!query-status!} 返回
\passthrough{\lstinline!running=true!}。测试结束后：

\begin{lstlisting}[language=sh]
qmp qmp_capabilities quit
\end{lstlisting}

QMP 的 quit 结束的是 QEMU 进程；日常开发关机也可以先在 guest 内正常
shutdown，再清理宿主运行目录。

\subsection{1.4
串口：按提示符同步，而不是盲目写入}\label{ux4e32ux53e3ux6309ux63d0ux793aux7b26ux540cux6b65ux800cux4e0dux662fux76f2ux76eeux5199ux5165}

串口自动化最容易犯的错误是一次写入大量命令。shell
尚未进入提示符时，后续输入可能拼接、丢失或落到错误程序中。

可靠的交互器应做到：

\begin{enumerate}
\def\labelenumi{\arabic{enumi}.}
\tightlist
\item
  等待 \passthrough{\lstinline!login:!}；
\item
  登录后等待明确的 shell prompt；
\item
  每发送一条命令，都等待下一条 prompt；
\item
  超时后保存最后一段输出并失败退出；
\item
  用唯一标记判断一组命令是否完成。
\end{enumerate}

基础冒烟命令可以包括：

\begin{lstlisting}[language=sh]
id
uname -a
freebsd-version
mount -p
cat -n /etc/fstab
ifconfig -l
jls
\end{lstlisting}

网络测试则先确认接口和地址，再测试网关或实际服务。没有配置
\passthrough{\lstinline!hostfwd!} 时，guest 内 sshd
正在运行并不意味着宿主可以直接 SSH。

\subsection{1.5
GDB：让暂停位置对应正确的内核}\label{gdbux8ba9ux6682ux505cux4f4dux7f6eux5bf9ux5e94ux6b63ux786eux7684ux5185ux6838}

只读检查：

\begin{lstlisting}[language=sh]
gdb -q -nx -batch \
    -ex 'set architecture i386:x86-64' \
    -ex "symbol-file $KERNEL_DEBUG" \
    -ex 'set remotetimeout 10' \
    -ex "target remote $RUNDIR/gdb.sock" \
    -ex 'info address mi_startup' \
    -ex 'info registers rip rax rbx' \
    -ex 'x/4i $rip' \
    -ex 'detach'
\end{lstlisting}

符号文件必须与运行内核匹配。连接 remote target 会暂停
vCPU；\passthrough{\lstinline!detach!}
才会让它继续。寄存器值随暂停位置变化，不应与历史数值逐字比较。

如果需要观察 MM 或 jail 路径，可以先从符号开始：

\begin{lstlisting}
info address vm_fault
info address vm_fault_cow
info address prison_deref
\end{lstlisting}

能解析符号只是起点，不等于已经证明某次执行经过该函数。后者还需要断点、条件、调用栈和
guest 侧触发之间的对应证据。

\subsection{1.6
可逆寄存器写入：验证通道，不制造结论}\label{ux53efux9006ux5bc4ux5b58ux5668ux5199ux5165ux9a8cux8bc1ux901aux9053ux4e0dux5236ux9020ux7ed3ux8bba}

下面实验对 vCPU0 的 RAX 做一次 bit 0 XOR，随后立即恢复：

\begin{lstlisting}[language=sh]
gdb -q -nx -batch \
    -ex 'set architecture i386:x86-64' \
    -ex 'set remotetimeout 10' \
    -ex "target remote $RUNDIR/gdb.sock" \
    -ex 'thread 1' \
    -ex 'set $saved_rax = $rax' \
    -ex 'set $injected_rax = $saved_rax ^ 0x1' \
    -ex 'set $rax = $injected_rax' \
    -ex 'set $injection_ok = ($rax == $injected_rax)' \
    -ex 'set $rax = $saved_rax' \
    -ex 'set $restore_ok = ($rax == $saved_rax)' \
    -ex 'printf "INJECTION_VERIFY=%d\n", $injection_ok' \
    -ex 'printf "RESTORE_VERIFY=%d\n", $restore_ok' \
    -ex 'detach'
\end{lstlisting}

两个标记都应为 1。之后再次查询 QMP，并在串口执行唯一标记：

\begin{lstlisting}[language=sh]
qmp qmp_capabilities query-status
# guest shell 中：
echo __POST_GDB_OK__
\end{lstlisting}

\textbf{已实测：} RAX 写入、读回、恢复和 VM 继续运行。

这个实验只证明 GDB stub
的读写恢复通道可用。它没有选择内核故障点，没有控制触发时机，也没有生成输入语料，因此不是
fuzzing。任意寄存器修改都可能破坏内核状态，只能在隔离、可丢弃的 VM
中进行。

\subsection{1.7 ATF 与
Kyua：测试程序和调度器各司其职}\label{atf-ux4e0e-kyuaux6d4bux8bd5ux7a0bux5e8fux548cux8c03ux5ea6ux5668ux5404ux53f8ux5176ux804c}

ATF 测试程序声明用例、元数据和断言；Kyua 读取
Kyuafile、隔离工作目录、调度执行并保存结果。

\begin{lstlisting}
test source
├── atf_test_case / atf_add_test_case
├── atf_set
├── atf_check
└── cleanup
        ↓
ATF runtime
        ↓
Kyua：调度、结果数据库和报告
\end{lstlisting}

常用操作：

\begin{lstlisting}[language=sh]
cd /usr/tests/usr.sbin/jail

# ATF 程序自身的列表
./jail_basic_test -l

# Kyua 合并 Kyuafile 后的有效元数据
kyua list -v jail_basic_test

# 整个测试程序或单个用例
kyua test jail_basic_test
kyua test jail_basic_test:basic

# 输出详细通过记录
kyua report --verbose --results-filter=passed

# 失败时重新执行并直接查看输出
kyua debug jail_basic_test:basic
\end{lstlisting}

直接执行 ATF 程序不会获得 Kyua 的隔离和结果管理。会修改
jail、接口或文件系统的测试，不应在未隔离的真实系统上直接运行。

\passthrough{\lstinline!atf\_check!} 默认要求退出码 0、stdout 和 stderr
为空。常用断言：

\begin{longtable}[]{@{}
  >{\raggedright\arraybackslash}p{(\linewidth - 2\tabcolsep) * \real{0.4800}}
  >{\raggedright\arraybackslash}p{(\linewidth - 2\tabcolsep) * \real{0.5200}}@{}}
\toprule\noalign{}
\begin{minipage}[b]{\linewidth}\raggedright
参数
\end{minipage} & \begin{minipage}[b]{\linewidth}\raggedright
含义
\end{minipage} \\
\midrule\noalign{}
\endhead
\bottomrule\noalign{}
\endlastfoot
\passthrough{\lstinline!-s exit:0!} /
\passthrough{\lstinline!-s not-exit:0!} & 精确退出码或要求非零 \\
\passthrough{\lstinline!-o empty!} / \passthrough{\lstinline!-e empty!}
& 输出必须为空 \\
\passthrough{\lstinline!-o match:"regexp"!} & 输出匹配正则 \\
\passthrough{\lstinline!-o save:file!} & 保存输出供后续检查 \\
\passthrough{\lstinline!-r timeout!} & 在期限内重试 \\
\end{longtable}

\subsection{1.8
jail\_basic\_test：结果和下一步同样重要}\label{jail_basic_testux7ed3ux679cux548cux4e0bux4e00ux6b65ux540cux6837ux91cdux8981}

\textbf{已实测：} 在本文构建的 FreeBSD guest
中，\passthrough{\lstinline!jail\_basic\_test!} 七个用例全部通过，Kyua
退出码为 0。

\begin{longtable}[]{@{}
  >{\raggedright\arraybackslash}p{(\linewidth - 2\tabcolsep) * \real{0.4800}}
  >{\raggedright\arraybackslash}p{(\linewidth - 2\tabcolsep) * \real{0.5200}}@{}}
\toprule\noalign{}
\begin{minipage}[b]{\linewidth}\raggedright
用例
\end{minipage} & \begin{minipage}[b]{\linewidth}\raggedright
覆盖重点
\end{minipage} \\
\midrule\noalign{}
\endhead
\bottomrule\noalign{}
\endlastfoot
\passthrough{\lstinline!basic!} & 创建、查询、修改、重启和删除 \\
\passthrough{\lstinline!commands!} & 配置文件及生命周期钩子 \\
\passthrough{\lstinline!jid\_name\_set!} & JID/name 的设置与展开 \\
\passthrough{\lstinline!list!} & VNET、epair 和接口列表 \\
\passthrough{\lstinline!nested!} & 层级 jail、children.max 与权限限制 \\
\passthrough{\lstinline!param\_consistency!} & 参数一致性和 poststop
环境 \\
\passthrough{\lstinline!setaudit!} &
\passthrough{\lstinline!allow.setaudit!} 的行为 \\
\end{longtable}

测试后的资源检查同样属于验收：

\begin{lstlisting}[language=sh]
jls
ifconfig -l
\end{lstlisting}

历史结果中没有残留 jail 或测试接口。verbose 报告还出现了
\passthrough{\lstinline!nestedjail not found!} 的 cleanup
信息；它没有使测试失败，却提供了一个值得研究的生命周期问题。\textbf{一条非致命
stderr 既不是通过的反证，也不是内核回归的证据。}

\textbf{未验证：} 完整 \passthrough{\lstinline!/usr/tests!}、所有 jail
配置组合以及 Kyua jail \passthrough{\lstinline!execenv!}。

到这里，验证环境已经具备两种互补能力：ATF/Kyua
给出可重复的功能结果，QEMU/GDB 给出实现层的观察窗口。后续 MM 和 jail
研究应同时保留这两类证据。

\section{2.trials: nixbsd on docker or
qemu+kvm}\label{trials-nixbsd-on-docker-or-qemukvm}

\subsection{2.1
试验的问题不是``能不能构建''，而是``能否自举开发环境''}\label{ux8bd5ux9a8cux7684ux95eeux9898ux4e0dux662fux80fdux4e0dux80fdux6784ux5efaux800cux662fux80fdux5426ux81eaux4e3eux5f00ux53d1ux73afux5883}

一个成功产物并不自动成为好用的开发环境。构建环境可能擅长固定依赖和缓存，却不擅长启动目标系统；虚拟机能够运行目标内核，却需要额外管理镜像、符号和控制台。

本章比较三条路线：

\begin{enumerate}
\def\labelenumi{\arabic{enumi}.}
\tightlist
\item
  Linux 容器中的 FreeBSD 交叉 builder；
\item
  Docker 中的 NixBSD builder 及其系统产物；
\item
  传统 QEMU/KVM guest 中的 FreeBSD 开发。
\end{enumerate}

重点放在环境边界和迭代成本。第 1 章已经给出的串口、QMP 和 GDB
操作不再重复。

\subsection{2.2 Linux builder
解决了什么}\label{linux-builder-ux89e3ux51b3ux4e86ux4ec0ux4e48}

FreeBSD 交叉 builder 把
Clang/LLD、bmake、自举工具和对象目录放进一个可重建环境。它适合：

\begin{itemize}
\tightlist
\item
  固定宿主工具链；
\item
  重复 world/kernel 构建；
\item
  保留对象和 release 缓存；
\item
  避免污染宿主全局环境。
\end{itemize}

但 Docker 容器共享宿主内核。Linux builder 中的 Bash、Nix 或 Clang 是
Linux 程序，生成的 FreeBSD \passthrough{\lstinline!/bin/sh!}
不会因为``构建成功''就能在同一个容器中运行。

因此 builder
的优势是\textbf{构建隔离}，不是\textbf{目标运行时}。把源码挂载进容器也不会改变这一点。

\subsection{2.3 NixBSD
builder：把声明式构建引入试验}\label{nixbsd-builderux628aux58f0ux660eux5f0fux6784ux5efaux5f15ux5165ux8bd5ux9a8c}

NixBSD 试验使用固定的上游提交：

\begin{lstlisting}
upstream: nix-community/nixbsd
commit:   fa821ab6c2a09dfb0411cbe52df25fe878b0f8b4
Nix:      2.31.2
\end{lstlisting}

该提交提供
\passthrough{\lstinline!.\#base.toplevel!}、\passthrough{\lstinline!.\#base.vm!}
等 flake 输出，但没有与本地试验相同的官方 Dockerfile、devShell 或 OCI
Live 镜像。这里的镜像是本地 builder 定义，不是上游发布物。

它的目标被刻意限定为：

\begin{itemize}
\tightlist
\item
  在 Linux 容器中提供 Nix；
\item
  固定 NixBSD 源码；
\item
  持久保存 \passthrough{\lstinline!/nix!} 和源码；
\item
  求值并构建 NixBSD 产物。
\end{itemize}

它从一开始就不是 NixBSD 用户态容器，也不负责启动 QEMU。

\subsection{2.4
镜像固定与入口约束}\label{ux955cux50cfux56faux5b9aux4e0eux5165ux53e3ux7ea6ux675f}

一个可信的 builder 镜像应满足：

\begin{enumerate}
\def\labelenumi{\arabic{enumi}.}
\tightlist
\item
  基础镜像固定版本和 digest，不使用 \passthrough{\lstinline!latest!}；
\item
  启用 \passthrough{\lstinline!nix-command flakes!}，不关闭 TLS、签名或
  substituter 校验；
\item
  构建时 clone 上游并 checkout 完整提交；
\item
  用 \passthrough{\lstinline!git rev-parse HEAD!}
  断言源码与预期完全一致；
\item
  默认工作目录为 \passthrough{\lstinline!/opt/nixbsd!}；
\item
  entrypoint 使用 \passthrough{\lstinline!exec "$@"!}
  转发命令，无参数时进入登录 shell。
\end{enumerate}

构建和检查可以概括为：

\begin{lstlisting}[language=sh]
export NIXBSD_IMAGE=nixbsd-live:fa821ab6
export NIXBSD_DOCKER="$WORK/trials/nixbsd-live"

docker build --pull=false -t "$NIXBSD_IMAGE" "$NIXBSD_DOCKER"

docker run --rm "$NIXBSD_IMAGE" bash -c '
    nix --version
    git --version
    git rev-parse HEAD
    nix config show experimental-features
'
\end{lstlisting}

这里的价值不在于容器标签，而在于\textbf{构建输入可解释}：镜像是哪一个版本、源码是哪一个提交、flake
特性是否开启，都可以在运行前检查。

\subsection{2.5
两个持久卷，两种生命周期}\label{ux4e24ux4e2aux6301ux4e45ux5377ux4e24ux79cdux751fux547dux5468ux671f}

试验中每次运行挂载两个命名卷：

\begin{lstlisting}[language=sh]
docker run --rm -it --name nixbsd-live \
    --mount source=nixbsd-nix,target=/nix \
    --mount source=nixbsd-src-fa821ab6,target=/opt/nixbsd \
    "$NIXBSD_IMAGE"
\end{lstlisting}

\begin{longtable}[]{@{}
  >{\raggedright\arraybackslash}p{(\linewidth - 2\tabcolsep) * \real{0.4800}}
  >{\raggedright\arraybackslash}p{(\linewidth - 2\tabcolsep) * \real{0.5200}}@{}}
\toprule\noalign{}
\begin{minipage}[b]{\linewidth}\raggedright
卷
\end{minipage} & \begin{minipage}[b]{\linewidth}\raggedright
保存内容
\end{minipage} \\
\midrule\noalign{}
\endhead
\bottomrule\noalign{}
\endlastfoot
\passthrough{\lstinline!nixbsd-nix!} & Nix 工具、flake 输入、构建缓存和
store \\
\passthrough{\lstinline!nixbsd-src-fa821ab6!} & 固定源码、后续修改和
result/GC root \\
\end{longtable}

这种拆分带来两个实际优势：

\begin{itemize}
\tightlist
\item
  删除临时容器不会删除缓存；
\item
  重建镜像不会让旧源码卷悄悄覆盖新提交。
\end{itemize}

源码卷按提交命名很重要。镜像升级而卷名不变时，旧
\passthrough{\lstinline!/opt/nixbsd!}
可能掩盖镜像中刚固定的源码。\passthrough{\lstinline!--rm!}
不等于删除命名卷，真正的风险来自显式 volume 清理。

\textbf{已实测：} 临时容器重建后，两个卷的内容和源码提交仍可复用。

\subsection{2.6 从求值到
dry-run：逐步扩大承诺}\label{ux4eceux6c42ux503cux5230-dry-runux9010ux6b65ux6269ux5927ux627fux8bfa}

进入 builder 后，先读取 flake，而不是直接进行大型构建：

\begin{lstlisting}[language=sh]
nix flake metadata . --no-write-lock-file

nix eval --no-write-lock-file \
    .#packages.x86_64-linux.base.toplevel.drvPath

nix build --dry-run --no-link --no-write-lock-file \
    .#base.toplevel
\end{lstlisting}

这三个步骤代表不同层次：

\begin{itemize}
\tightlist
\item
  metadata 证明 flake 输入和源码版本可读取；
\item
  \passthrough{\lstinline!drvPath!} 证明目标可以求值；
\item
  dry-run 证明 Nix 能生成构建计划。
\end{itemize}

dry-run 不会完成计划中的编译和下载，因此不能写成``系统已构建''。如果
flake 建议额外 substituter 或公钥，应先审查固定提交的
\passthrough{\lstinline!nixConfig!}，不要盲目使用
\passthrough{\lstinline!--accept-flake-config!}。

\textbf{已实测：} 上述 metadata、目标求值和 dry-run 均成功。

\subsection{2.7 实际编译
toplevel}\label{ux5b9eux9645ux7f16ux8bd1-toplevel}

真正跨过``只有计划''边界的是：

\begin{lstlisting}[language=sh]
nix build --no-write-lock-file \
    --max-jobs 4 --cores 4 -L \
    .#base.toplevel
\end{lstlisting}

\textbf{已实测：} \passthrough{\lstinline!.\#base.toplevel!}
完整构建成功，输出目标为
\passthrough{\lstinline!x86\_64-freebsd!}，并检查了以下结构：

\begin{lstlisting}
toplevel/
├── activate
├── init
├── kernel
├── etc
└── sw
\end{lstlisting}

\passthrough{\lstinline!activate!} 和 \passthrough{\lstinline!init!}
可执行，\passthrough{\lstinline!kernel!}
链接目标存在且非空；\passthrough{\lstinline!init!} 来自 FreeBSD
\passthrough{\lstinline!init-static!}，\passthrough{\lstinline!kernel!}
来自带模块的内核输出。随后通过普通 \passthrough{\lstinline!nix build!}
建立持久 result/GC root，降低构建结果被垃圾回收的风险：

\begin{lstlisting}[language=sh]
nix build --no-write-lock-file .#base.toplevel
\end{lstlisting}

这个 result 位于源码命名卷中。它证明一套 NixBSD
系统闭包已经形成，但仍不证明该闭包已经启动。

构建中断时，应使用同一组持久卷续建并以最终退出码和产物检查为准，不能根据日志末尾停在哪里猜测成功。

\subsection{2.8 toplevel
之后，仍然隔着运行时}\label{toplevel-ux4e4bux540eux4ecdux7136ux9694ux7740ux8fd0ux884cux65f6}

\begin{longtable}[]{@{}
  >{\raggedright\arraybackslash}p{(\linewidth - 4\tabcolsep) * \real{0.2400}}
  >{\raggedright\arraybackslash}p{(\linewidth - 4\tabcolsep) * \real{0.4000}}
  >{\raggedright\arraybackslash}p{(\linewidth - 4\tabcolsep) * \real{0.3600}}@{}}
\toprule\noalign{}
\begin{minipage}[b]{\linewidth}\raggedright
目标
\end{minipage} & \begin{minipage}[b]{\linewidth}\raggedright
本文状态
\end{minipage} & \begin{minipage}[b]{\linewidth}\raggedright
能得出的结论
\end{minipage} \\
\midrule\noalign{}
\endhead
\bottomrule\noalign{}
\endlastfoot
\passthrough{\lstinline!.\#base.toplevel!} & 已实测构建 &
系统闭包可构建、结构可检查 \\
\passthrough{\lstinline!.\#base.vm!} & 仅源码审查 & 上游存在 VM runner
路径 \\
\passthrough{\lstinline!.\#base.systemImage!} & 仅源码审查 & 上游定义了
UEFI qcow2 生成逻辑 \\
NixBSD 用户态执行 & 未验证 & 不能由 toplevel 推断 \\
NixBSD ISO & 未提供 & 不能套用 FreeBSD disc1 安装流程 \\
NixBSD 物理安装 & 未验证 & 不具备实测安装结论 \\
\end{longtable}

源码审查显示，VM 路线涉及
\passthrough{\lstinline!vmClosureInfo!}、OVMF、可写 CoW 磁盘和 EFI
变量；\passthrough{\lstinline!systemImage!} 是 GPT/ESP/UFS 的
qcow2，不是 ISO。上游还提示 VM 闭包可能消耗大量磁盘。

这些资料足以解释下一步可能是什么，却不足以写成操作结果。尤其是通用
\passthrough{\lstinline!nixos-install!} 文档中的 NixOS/ext4/GRUB
示例，不能直接冒充 NixBSD/FreeBSD 物理安装流程。

\subsection{2.9 容器内 FreeBSD qcow2
安装}\label{ux5bb9ux5668ux5185-freebsd-qcow2-ux5b89ux88c5}

\textbf{状态：暂停，留空。}

容器封装 FreeBSD VM
的安装盘启动、完整安装和安装后持久启动没有完成。本章不补充安装成功结论，也不把
Live ISO 的启动记录当作安装记录。

\subsection{2.10 最终选择：传统
KVM+QEMU+GDB}\label{ux6700ux7ec8ux9009ux62e9ux4f20ux7edf-kvmqemugdb}

三条路线的优势并不相同：

\begin{longtable}[]{@{}
  >{\raggedright\arraybackslash}p{(\linewidth - 4\tabcolsep) * \real{0.3333}}
  >{\raggedright\arraybackslash}p{(\linewidth - 4\tabcolsep) * \real{0.3333}}
  >{\raggedright\arraybackslash}p{(\linewidth - 4\tabcolsep) * \real{0.3333}}@{}}
\toprule\noalign{}
\begin{minipage}[b]{\linewidth}\raggedright
路线
\end{minipage} & \begin{minipage}[b]{\linewidth}\raggedright
最强之处
\end{minipage} & \begin{minipage}[b]{\linewidth}\raggedright
主要代价
\end{minipage} \\
\midrule\noalign{}
\endhead
\bottomrule\noalign{}
\endlastfoot
Linux FreeBSD builder & 构建依赖固定、缓存稳定 & 不提供 FreeBSD
runtime \\
Docker NixBSD builder & 声明式闭包、源码与 store 可持久化 &
构建与启动之间仍有较长链路 \\
KVM+QEMU+GDB & 直接控制内核、控制台和调试接口 &
需要自行管理镜像、符号和运行目录 \\
\end{longtable}

对于 MM 和 jail 研究，主要迭代不是``再生成一次声明式系统''，而是：

\begin{lstlisting}
修改内核 → 启动 → 触发路径 → 观察 → 修改
\end{lstlisting}

传统虚拟化让这条链路更短。KVM 可用时还能减少 TCG
的执行开销；即使暂时没有 KVM，同一套 QEMU 控制面也可以先用 TCG
完成验证。

因此后续路线确定为：

\begin{quote}
\textbf{放弃继续推进 NixBSD，把 Docker/Nix 留在构建试验记录中；MM 与
jail 开发回到传统 KVM+QEMU+GDB。}
\end{quote}

这是一项路线选择，不是对 NixBSD 项目价值的否定，也不表示当前宿主已经验证
KVM。

\section{3.real machine deployment}\label{real-machine-deployment}

\subsection{3.1
先识别介质格式}\label{ux5148ux8bc6ux522bux4ecbux8d28ux683cux5f0f}

把虚拟机产物带到物理机之前，必须先确认自己拿的是什么：

\begin{longtable}[]{@{}
  >{\raggedright\arraybackslash}p{(\linewidth - 4\tabcolsep) * \real{0.2400}}
  >{\raggedright\arraybackslash}p{(\linewidth - 4\tabcolsep) * \real{0.4000}}
  >{\raggedright\arraybackslash}p{(\linewidth - 4\tabcolsep) * \real{0.3600}}@{}}
\toprule\noalign{}
\begin{minipage}[b]{\linewidth}\raggedright
格式
\end{minipage} & \begin{minipage}[b]{\linewidth}\raggedright
典型用途
\end{minipage} & \begin{minipage}[b]{\linewidth}\raggedright
不能混淆的地方
\end{minipage} \\
\midrule\noalign{}
\endhead
\bottomrule\noalign{}
\endlastfoot
ISO & 光盘/安装或 Live 介质 & 与 qcow2 的用途不同 \\
raw & 原始块设备内容 & 写入前必须确认完整磁盘布局 \\
qcow2 & QEMU 写时复制磁盘 & 不是可直接等同 ISO 的安装介质 \\
\end{longtable}

转换格式只改变容器，不会自动增加固件支持。BIOS/UEFI、控制台、磁盘驱动和启动项仍需分别验证。

\subsection{3.2
控制台必须适合目标机器}\label{ux63a7ux5236ux53f0ux5fc5ux987bux9002ux5408ux76eeux6807ux673aux5668}

为 QEMU \passthrough{\lstinline!-serial stdio!} 制作的 Live ISO
可能包含：

\begin{lstlisting}[language=sh]
console="comconsole"
comconsole_speed="115200"
boot_serial="YES"
\end{lstlisting}

这些设置适合串口自动化，却可能让物理显示器在后续启动阶段看不到安装器或
login。用于真实机器前，应根据硬件选择视频控制台或确认串口设备存在，不能把
QEMU 的可见性直接外推到物理机。

\subsection{3.3 制作
USB：先确认整块设备，再写入}\label{ux5236ux4f5c-usbux5148ux786eux8ba4ux6574ux5757ux8bbeux5907ux518dux5199ux5165}

设置介质与目标变量：

\begin{lstlisting}[language=sh]
export LIVE_ISO="$IMAGES/freebsd-live-amd64.iso"
export USB_DISK=/dev/CHANGE_ME
\end{lstlisting}

写入前至少完成以下人工检查：

\begin{enumerate}
\def\labelenumi{\arabic{enumi}.}
\tightlist
\item
  \passthrough{\lstinline!LIVE\_ISO!} 是预期文件并记录 SHA-256；
\item
  \passthrough{\lstinline!USB\_DISK!}
  是专门的、可整体清除的可移动块设备；
\item
  目标是整块磁盘，不是其中一个分区；
\item
  所有挂载自该设备的分区都已卸载；
\item
  已再次确认内部系统盘不在目标中。
\end{enumerate}

可用 \passthrough{\lstinline!lsblk!}
的型号、容量、传输类型和挂载点交叉确认，不要只依赖设备字母顺序。

写入与同步：

\begin{lstlisting}[language=sh]
sudo dd if="$LIVE_ISO" of="$USB_DISK" \
    bs=4M status=progress oflag=sync conv=fsync
sync
\end{lstlisting}

按 ISO 长度逐字节比较，避免 U 盘尾部额外空间影响判断：

\begin{lstlisting}[language=sh]
ISO_SIZE="$(stat -c %s "$LIVE_ISO")"
sudo cmp -n "$ISO_SIZE" "$LIVE_ISO" "$USB_DISK"
\end{lstlisting}

\passthrough{\lstinline!cmp!} 无差异输出才支持``已写入内容与 ISO
一致''。它仍不证明固件能够启动，也不证明安装到内部磁盘后能够运行。

\subsection{3.4
当前证据边界}\label{ux5f53ux524dux8bc1ux636eux8fb9ux754c}

\textbf{已实测：} FreeBSD Live ISO 写入专用 SanDisk U
盘的整块设备，完成同步、逐字节比较和启动结构检查。

\textbf{未验证：}

\begin{itemize}
\tightlist
\item
  物理机从该介质完整启动并进入 Live；
\item
  物理机硬件发现和网络；
\item
  安装到内部磁盘；
\item
  安装后的首次启动与持久使用。
\end{itemize}

近期围绕已有 ZFS 启动环境的排查，也不能替代上述 Live
介质验证。真机部署仍应把``介质制作成功''和``机器启动成功''分成两个验收点。

\section{4.kernel eject and fuzz}\label{kernel-eject-and-fuzz}

\textbf{状态：尚未开始。}

后续范围只包括 FreeBSD \textbf{MM} 与 \textbf{jail}。NixBSD
不再进入研究队列，其他内核子系统也不通过全量 stress 间接带入。

开始之前需要建立以下门槛：

\begin{enumerate}
\def\labelenumi{\arabic{enumi}.}
\tightlist
\item
  \textbf{隔离：} 所有 stress/fuzz 只在专用 FreeBSD guest 内执行，基础
  qcow2 使用 \passthrough{\lstinline!-snapshot!}。
\item
  \textbf{范围：} stress2 只能逐脚本审查后进入 MM/jail
  allowlist，不能运行全量套件。
\item
  \textbf{边界：} 不连接宿主块设备、PCI 设备或可写共享；fuzz 默认
  \passthrough{\lstinline!-nic none!}。
\item
  \textbf{时限：} 每轮有单命令和总超时，保留 serial/QMP/GDB
  证据，并能独立清理。
\item
  \textbf{插桩：} GENERIC 不应被假定具有 KCOV；需要单独验证 KASAN/KCOV
  内核和 \passthrough{\lstinline!/dev/kcov!}。
\item
  \textbf{结论：} 超时、panic、测试失败、cleanup
  噪声和预期异常必须分类记录，不能混为``发现漏洞''。
\end{enumerate}

第 1 章的可逆 RAX 实验只是 GDB 通道验证。真正的内核故障注入或 fuzz
还需要明确触发点、输入生成、覆盖反馈、状态恢复和结果分类。

在这些条件满足之前，本章保留为空白，不预先填写任何 crash、coverage
或漏洞结论。

\end{document}
